% Une ligne commentaire débute par le caractère 

\documentclass[a4paper]{article}

% Options possibles : 10pt, 11pt, 12pt (taille de la fonte)
%                     oneside, twoside (recto simple, recto-verso)
%                     draft, final (stade de développement)

\usepackage[utf8]{inputenc}   % LaTeX, comprends les accents !
\usepackage[T1]{fontenc}      % Police contenant les caractères français
\usepackage[francais]{babel}  
\usepackage{fullpage}
\usepackage[dvipsnames]{xcolor}



\usepackage{graphicx}  % pour inclure des images

%\pagestyle{headings}        % Pour mettre des entêtes avec les titres
                              % des sections en haut de page

 \title{  Simulation de fluides en 2D\\         % Les paramètres du titre : titre, auteur, date
  Projet de programmation 2}          
\author{ PEYREROL Axel, MILEZI Gaëlle,\\ BARRAUX Satine,  DELRIEU Saurane\\
  \emph{L3 informatique}\\
  Faculté des Sciences\\
Université de Montpellier.}
\date{2025-2026}             



\begin{document}

\maketitle                    % Faire un titre utilisant les données
                              % passées à \title, \author et \date

%\begin{center}               % pour centrer 
 % \includegraphics[scale=1]{}   % insertion d'une image
%\end{center}

\begin{abstract}     % Résumé du travail

  \emph{Description très succinte du problème et des différentes étapes de réalisation}

\end{abstract}
\newpage
%\dominitoc  % initializer les minitoc
\tableofcontents
\newpage
\section*{Introduction} % Commencer une section, * permet de ne pas numéroté l'introduction
\emph{1-2 pages}
	%Présentation du Sujet (1-2 pages) 
\begin{itemize}%pour faire une liste à puces
\item Présenter le problème étudié et le contexte dans lequel le projet se positionne.
\begin{itemize}
\item Motiver l'intérêt du problème étudié par rapport à votre parcours d'études et au monde de l'informatique.
\item	Présenter les différentes approches possibles pour la résolution du problème, et en particulier celle choisie.
\item Donner le cahier des charges détaillé
\end{itemize}

\section{Technologies utilisées}
\emph{1-2 pages}
\begin{itemize}
\item Langages et bibliothèques : présenter les langages de programmation et les outils utilisés
\item Intéret de ces choix : justifier l'intéret et les choix de ces langages 
\end{itemize}


\section{Conception, Modélisation, Implémentation} 
\emph{8 à 15 pages}

\begin{enumerate}%pour faire une liste numérotée
\item Présenter les développements réalisés : une fenêtre, une zone de rendu et un panneau de paramètres.
\item Présenter les principaux modules du projet. (Utiliser le langage UML pour la modélisation : donner le diagramme de cas d'utilisation et le diagramme des classes.)
\item Décrire les fonctionnalités de l'interface graphique implémentée (champ de simulatio, boutons, affichage de l'obstacle etc...).
    \begin{itemize}
    {\color{WildStrawberry}
    \item  obstacle: dans un premiers temps, il n'y a eu q'un seul obstacle avec un bouton on/of sur celui ci. Ensuite ce bouton fut remplacer par un bouton +ou- afin de pouvoir plusieur obstacle a la fois. De plus , il est egalement possible de mettre des obsatcle en forme de etoile a 6 branches (2 triangle equilaterale imbriquer + probleme de bord) et de mettre des coeurs ( tres bizarre). de plus l'ensemble de ces obstacles peuvent etre deplacer en cliquant dessus. le premier afficher est le premeier deplacer si ils sont tous les uns sur les autres car c'est le premeires obstaclede la liste. de plus le nombres d'obstacles par forme est afficher.
    }
    \end{itemize}
    
\item Décrire le format des données en entrée ou encore les conventions utilisées pour les entrées de vos programmes. Décrire les procédures de lecture et validation des entrées. (représentation du fluide, grille, stockage des paramètres)
\item Statistiques : nombre de modules/composantes/classes/scripts développés. Nombre de lignes de code.
\end{enumerate}

\section{Algorithmes et Structures de Données}
\emph{2 à 3 pages}
\begin{enumerate}%pour faire une liste numérotée
\item 	Présenter les principales structures de données définies dans le cadre du projet.
\item 	Présenter les principaux algorithmes implémentés. En illustrer le fonctionnement avec des exemples. Attention : il s'agit ici de choisir 1 ou 2 algorithmes intéressants, et non pas de présenter tous les algorithmes implémentés.
\item 	Évaluer la complexité théorique en temps des algorithmes présentés
\end{enumerate}

\section{Analyse des résultats }
\emph{2 à 4 pages}
L'analyse expérimentale sera faites dans cette section.   Elle sera plus ou moins importante dans votre projet comme pour les deux points précédent, cette section devra donc être développée en conséquence.

\begin{enumerate}%pour faire une liste numérotée
\item 	Illustrer les performances ainsi que l'efficacité du logiciel implémenté à l'aide de graphiques.
\item Analyser (et comparer, si plusieurs) les performances des solutions implémentées.
\item	Présenter les bancs d'essais (ou les procédures utilisées pour la génération des données) utilisés pour les tests du logiciel.
\end{enumerate}

\section{Gestion du Projet}
\emph{ (1-2 pages)}
\begin{enumerate}%pour faire une liste numérotée
\item 	Présenter la gestion du projet et les documents de planification rédigés (par exemple, le diagramme de Gantt).
    \begin{itemize}
    {\color{WildStrawberry}
    \item  mise en place du project par tout le monde (generation du fluide)
    \item partage des taches: Satine sur les obstacle (creation d'obstacle en forme de boule, d'etoile ou de coeur, mise en place de la possibilite d'en avoir plusieur en me e temps), Saurane sur le mélange de couleur, Gaelle et Axel sur les curseur, Gaelle sur la mise en forme, Axel sur le bluetooth.
    }
    \end{itemize}

\item	Discuter les changements majeurs effectués en cours de projet. 
\begin{itemize}
    {\color{WildStrawberry}
    \item  premiere implementation afin de ce familiariser avec le ave le C++ avant de passer sur le texte de Donald E. Knuth et ainsi implementer ce code 
    \item probleme de dependance sur les obstacle donc modification sur le main afin de proceder en liste. Probleme de bort :creation de fonction d'affichage et de fonction de bord.
    }
\end{itemize}

    
\end{enumerate}

\section{Bilan et Conclusions}

	Indiquer les fonctionnalités mises en œuvre par rapport au cahier des charges de départ, les points ouvertes.
	Les perspectives du projet décrivent ce que l’on pourrait ajouter, modifier, réécrire afin d’améliorer le projet.
    \begin{itemize}
    {\color{WildStrawberry}
    \item  fluide avec curseur et obstacle et bluetooth
    }
    \end{itemize}
	
\begin{thebibliography}{9}
\bibitem{texbook}
Donald E. Knuth (1986) \emph{The \TeX{} Book}, Addison-Wesley Professional.

\bibitem{lamport94}
Leslie Lamport (1994) \emph{\LaTeX: a document preparation system}, Addison
Wesley, Massachusetts, 2nd ed.
\end{thebibliography}



\end{document}